\section{模型与推断目标}
\label{sec:background}

\subsection{观测尺度已知的分枝实验}
\label{sec:bg-model}

考虑由 $n\ge1$ 例原发感染构成的一代传播。给定再生数 $R_t>0$，原发病例 $i$ 独立产生负二项后代\citep{lloydsmith2005superspreading}，
\begin{equation}
  X_{i,t}\mid R_t\sim\NB\!\left(k,\frac{k}{k+R_t}\right),\qquad
  \E(X_{i,t}\mid R_t)=R_t,\quad
  \Var(X_{i,t}\mid R_t)=R_t+\frac{R_t^2}{k}.
  \label{eq:offspring}
\end{equation}
其中 $k>0$ 控制个体异质性。每例继发感染以概率 $\rho\in(0,1]$ 独立报告。求和与二项稀疏化给出
\begin{equation}
\begin{aligned}
  C_t\mid n,R_t&\sim\NB\!\left(nk,\frac{k}{k+\rho R_t}\right),\\
  \E(C_t\mid n,R_t)&=\rho nR_t,\qquad
  \Var(C_t\mid n,R_t)=\rho nR_t+\frac{\rho^2nR_t^2}{k}.
\end{aligned}
  \label{eq:observed}
\end{equation}
定义 $m=\rho n$ 和 $\Rhat_t=C_t/m$。设计规模 $m$ 是 $R_t=1$ 时的期望报告数，而非实现计数 $C_t$。

基准假设为：$\rho$ 固定已知且报告相互独立（H1）；后代在条件下独立并服从式~(\ref{eq:offspring})（H2）；$k$ 固定，且在似然计算中已知（H3）；$n$ 固定或外生给定，并与传播独立（H4）。结果以所选设计为条件，各期具有相同的 $m$。记 $\E(R_t)=\Rbar>0$、$\Var(R_t)=v_R<\infty$。对于重复时期，另假定给定完整传播轨迹后各期计数独立，且条件分布为式~(\ref{eq:observed})。这里描述的是重复分枝实验，并不主张自然演化的发病数具有外生性。

观测尺度至关重要。未知报告比例的稳定性不足以确立无偏性：若除以假定的 $\rho_0$，条件均值为 $(\rho/\rho_0)R_t$，增加病例无法消除相应的乘性偏差。由含噪观测发病数构造的比值也具有随机分母，不必具备 $C_t/(\rho n)$ 的性质。因此，基准模型是一种明确的参照实验，而不是适用于所有监测统计量的直接模型。

条件计数方差同时包含后代随机性与独立漏报。先以继发感染总数为条件，可将两者贡献分别写为 $(R_t+R_t^2/k)/n$ 和 $(1-\rho)R_t/(\rho n)$，二者之和等于上述方差。特别地，与 $R_t$ 成线性的项并不全是报告噪声：完全报告时它仍然存在。这一区分避免将纯代数分项误认为唯一的生物学归因。

\subsection{当期值、实现均值与过程均值}
\label{sec:bg-variance}

对于当期目标，$\E(\Rhat_t\mid R_t)=R_t$，条件方差为 $[R_t+\rho R_t^2/k]/m$，达到单参数计数模型中的条件 Cram\'er--Rao 下界（附录~\ref{app:proofs}）。该陈述针对在此似然下对每个 $R_t$ 都无偏的估计量，不是普适的极小极大下界，也不涵盖有偏、分层或使用先验信息的方法。

\begin{proposition}[单个时期的两个目标]
\label{prop:average}
在 H1--H4 下，令 $A=\Rbar+\rho(\Rbar^2+v_R)/k$，则
\begin{equation}
\begin{aligned}
  \E[(\Rhat_t-R_t)^2]&=\frac{A}{m},\\
  V(m):=\E[(\Rhat_t-\Rbar)^2]&=v_R+\frac{A}{m}.
\end{aligned}
  \label{eq:var-average}
\end{equation}
在仅基于当期计数且对 $R_t$ 条件无偏的估计量类中，$\Rhat_t$ 同时最小化这两种风险。
\end{proposition}

第一种风险按指定的 $R_t$ 分布对条件估计误差取平均，既不是每个固定取值下的条件风险，也不是某一时期的后验可信区间。第二种风险以不同目标评价同一个观测。单期比较只涉及边际矩，因此不要求传播时期彼此独立；跨期汇总风险时，独立性才发挥作用。

这一区别来自
\begin{equation}
  \Rhat_t-\Rbar=(\Rhat_t-R_t)+(R_t-\Rbar).
  \label{eq:error-split}
\end{equation}
条件无偏性使两项不相关。跨期传播变异属于当期信号，却为过程均值贡献抽样误差。后者是指定平稳分布下的均值，不会自动等同于非平稳流行过程的季节平均。两项贡献具有整群抽样中熟悉的单元内与单元间结构\citep{cochran1977sampling, engen1998demographic}。

当 $v_R>0$ 时，定义平衡规模 $\mx=A/v_R$。固定 $\rho$、$k$ 与传播分布，通过 $n$ 增加 $m$ 不改变 $A$；在 $m=\mx$ 处，$V(m)$ 的两项贡献相等。这一规模描述所选估计量与目标，并不意味着所有信息的丧失；改变 $\rho$ 则会同时改变 $A$ 和 $m$。覆盖与预算计算仅作为符号化扩展列于附录~\ref{app:extensions}。

\begin{corollary}[固定的一组已观测时期]
\label{cor:targets}
在 H1--H4 及给定 $R_{1:T}$ 后计数条件独立的假设下，定义 $\Rbar_T=T^{-1}\sum_{t=1}^T R_t$ 和 $\widehat{\Rbar}_T=T^{-1}\sum_{t=1}^T\Rhat_t$。对整数 $T\ge1$，若 $R_t$ 具有共同边际矩，则
\begin{equation}
\begin{aligned}
  \E[(\widehat{\Rbar}_T-\Rbar_T)^2]&=\frac{A}{Tm},\\
  \E[(\widehat{\Rbar}_T-\Rbar)^2]&=B_T+\frac{A}{Tm},
  \qquad B_T=\Var(\Rbar_T).
\end{aligned}
  \label{eq:target-means}
\end{equation}
若 $R_t$ 独立同分布，则 $B_T=v_R/T$。对满足 $\Cov(R_t,R_{t+h})=v_R\phi^{|h|}$、$|\phi|<1$ 的平稳过程，有限 $T$ 下的精确表达式为 $B_T=v_RS(\phi,T)/T^2$，其中
\begin{equation}
  S(\phi,T)=T+2\sum_{h=1}^{T-1}(T-h)\phi^h
  =\frac{T(1+\phi)}{1-\phi}-\frac{2\phi(1-\phi^T)}{(1-\phi)^2}.
  \label{eq:ar-sum}
\end{equation}
\end{corollary}

实现均值的恒等式不要求各期传播水平独立，只要求计数误差条件独立且具有所述共同矩。固定 $T$ 而增加病例，可以消除实现均值的计数误差；过程均值的风险则趋于 $B_T$。追加时期应对的是这一项，其贡献取决于相依性。预测未来值是另一个目标，见附录~\ref{app:extensions}。

平稳协方差公式使用整数时期和完整的有限求和，而非长序列近似。正的持续性通常减小相邻时期带来的收益，负自相关则可能增加收益。不能自动将序列相关过程或均值变化的流行季替换为独立抽取。对于非平稳传播，实现均值仍有明确定义，但共同过程均值的解释需要额外模型。

空间平均具有类似条件。若 $J$ 个单元的传播方差均为 $v_R$、平均两两相关系数为 $\bar c$，其等权传播均值的方差为
\begin{equation}
  \Var\!\left(J^{-1}\sum_{j=1}^J R_{j,t}\right)
  =v_R\left[\bar c+\frac{1-\bar c}{J}\right].
  \label{eq:pooling}
\end{equation}
这是协方差恒等式，并非经验回归截距的分解。方差不等、权重不等或相关结构变化时，需使用对应的协方差矩阵。

\subsection{不可观测的报告波动与识别}
\label{sec:bg-reporting}

将 H1 替换为共同的逐期报告概率 $p_t\in[0,1]$：其独立于 $R_t$ 及设计，已知均值为 $\rho>0$，且 $\Var(p_t)=\rho^2\tau^2$（H5）。给定 $p_t$ 时，各例感染仍独立报告。必有 $\tau^2\le(1-\rho)/\rho$。估计量仍除以 $m=\rho n$，并不观测 $p_t$。

\begin{proposition}[报告波动与矩不可识别性]
\label{prop:report-noise}
在 H2--H5 下，记 $Q=\Rbar^2+v_R$，则
\begin{equation}
\begin{aligned}
  \E(\Rhat_t)&=\Rbar,\qquad
  V_H(m):=\Var(\Rhat_t)=v_R+v_H+\frac{A_H}{m},\\
  v_H&=\tau^2Q,\qquad
  A_H=\Rbar+\frac{\rho(1+\tau^2)Q}{k}.
\end{aligned}
  \label{eq:vH-AH}
\end{equation}
当 $m$、$\rho$ 和 $k$ 已知时，计数的前两个矩仅通过 $U=(1+\tau^2)(\Rbar^2+v_R)$ 依赖 $v_R$ 与 $\tau^2$，因此一般不能分别识别传播方差与报告波动。
\end{proposition}

同时给定 $R_t$ 与 $p_t$ 时，估计量的均值为 $(p_t/\rho)R_t$。只有对均值正确且独立的报告过程取平均后，才恢复对 $R_t$ 的无偏性。这说明为何同一时期内增加感染数无法平均掉共同报告冲击。如果能测量当期报告概率且它为正，正确调整的估计量将对应另一个问题；测量该概率的不确定性本身也需要传播。

这是指定模型族中的矩不可识别结果，并不意味着外部报告测量、附加假设下的时间信息，或限制更强的完整分布永远无法识别机制。边界情形也可能带来进一步约束。报告与传播独立是一项实质性条件；若疾病负担改变报告实践，上述分解未必成立。

\begin{corollary}[报告波动下的当期值与实现均值风险]
\label{cor:report-current}
在 H2--H5 下，
\begin{equation}
  \E[(\Rhat_t-R_t)^2]=v_H+\frac{A_H}{m}.
  \label{eq:vH-current}
\end{equation}
若另外假定各期 $p_t$ 独立同分布、独立于完整传播轨迹，且给定两条轨迹后各期计数独立，则
\begin{equation}
  \E[(\widehat{\Rbar}_T-\Rbar_T)^2]=\frac{v_H+A_H/m}{T},\qquad
  \E[(\widehat{\Rbar}_T-\Rbar)^2]=B_T+\frac{v_H+A_H/m}{T}.
  \label{eq:report-means}
\end{equation}
\end{corollary}

因此，实现目标中没有传播变异项，不代表没有报告误差。固定 $T$ 时，随着 $m$ 增加，不可观测的报告波动仍然存在。报告的序列相依性会增加协方差项，稳定但未知的平均报告比例也会引入偏差。表~\ref{tab:targets} 汇总了各恒等式及其条件。

\begin{table}[t]
\centering
\small
\setlength{\tabcolsep}{4pt}
\caption{推断目标及精确均方误差。固定报告列假设 H1--H4；波动报告列假设 H2--H5，且对跨期均值另要求各期报告概率独立、独立于完整传播轨迹。各期计数误差条件独立。$B_T=\Var(\Rbar_T)$；传播独立同分布时等于 $v_R/T$，具有所述平稳自协方差时等于 $v_RS(\phi,T)/T^2$。}
\label{tab:targets}
\begin{tabularx}{\linewidth}{>{\raggedright\arraybackslash}X>{\centering\arraybackslash}p{0.16\linewidth}>{\centering\arraybackslash}p{0.24\linewidth}>{\centering\arraybackslash}p{0.27\linewidth}}
\toprule
目标 & 估计量 & 固定且已知的报告 & 不可观测的波动报告 \\
\midrule
当期值 $R_t$ & $\Rhat_t$ & $A/m$ & $v_H+A_H/m$ \\
实现均值 $\Rbar_T$ & $\widehat{\Rbar}_T$ & $A/(Tm)$ & $(v_H+A_H/m)/T$ \\
过程均值 $\Rbar$ & $\widehat{\Rbar}_T$ & $B_T+A/(Tm)$ & $B_T+(v_H+A_H/m)/T$ \\
\bottomrule
\end{tabularx}
\end{table}

\subsection{内生计数规模与表观统计量}
\label{sec:bg-scope}

若 $m=\rho n$ 随机且与传播关联，即使保留给定 $(R_t,m)$ 的条件计数分布与无偏性，也只能得到
\begin{equation}
\begin{aligned}
 \E[(\Rhat_t-\Rbar)^2\mid m]
 &=\frac{\E(R_t\mid m)+\rho\E(R_t^2\mid m)/k}{m}\\
 &\quad+\Var(R_t\mid m)+[\E(R_t\mid m)-\Rbar]^2.
\end{aligned}
\label{eq:endogenous-risk}
\end{equation}
这些条件矩可能随规模改变，因此推不出常数截距加常数系数乘计数倒数的形式。以观测到的结果计数为条件，限制更强：这种条件化未必保留原先假定的条件观测分布。仅采用更新方程表述并不能解决上述任何问题。

因此，实证响应定义为表观统计量相对于指定中心的相对偏差平方（第~\ref{sec:meth-stats}~节），而非相对于已知 $R_t$ 的误差。其仿射截距和斜率不必等于 $v_R/\Rbar^2$ 与 $A/\Rbar^2$。延迟、聚合、分母噪声与变化的条件均值都会影响这一映射。校准规定的是依赖模型的联系（第~\ref{sec:meth-inversion}~节），不会把经验截距变成生物学来源的精确加性分解。
